\documentclass[twocolumn]{aastex631}

\usepackage{amsmath}
\usepackage{multirow}
\usepackage{natbib}
\usepackage{graphicx} 
\usepackage{aas_macros}

\begin{document}

Degenerate Higher-Order Scalar-Tensor (DHOST) theories have emerged as a compelling class of modifications to General Relativity, offering a rich theoretical framework to address fundamental questions in cosmology, such as the mechanism behind the late-time cosmic acceleration. By construction, these theories are classically stable, avoiding the ghost-like instabilities that often plague theories with higher-order derivatives. However, for any theory to be considered a viable candidate for describing nature, it must not only be classically well-behaved but also possess a consistent quantum description. The path to quantizing DHOST theories is obstructed by a subtle but profound feature: their gauge symmetries are field-dependent. This means that the parameters governing the transformations depend on the dynamical fields of the theory itself, a characteristic that complicates the underlying algebraic structure and poses a significant challenge to standard quantization protocols.This paper contends that the difficulties in quantizing DHOST theories stem from a foundational misunderstanding of their symmetry structure. We posit that the field-dependent gauge symmetries do not form a closed Lie algebra, as is implicitly assumed in conventional approaches. Instead, we argue they form an open algebra, a structure where the commutator of two infinitesimal gauge transformations fails to close off-shell. Specifically, the commutator yields another gauge transformation only after invoking the classical equations of motion, taking the schematic form $[\delta_{\epsilon_1}, \delta_{\epsilon_2}] \sim \delta_{\epsilon_3} + (\text{EOM}) \times (\text{coefficients})$. The presence of these on-shell terms renders standard quantization techniques, such as the Faddeev-Popov procedure, inapplicable, as they are predicated on a strictly off-shell algebraic closure. The appropriate and necessary framework for quantizing theories with open gauge algebras is the Batalin-Vilkovisky (BV) formalism, a powerful generalization of BRST quantization designed precisely for such intricate gauge systems.Our investigation proceeds in two main stages. First, we establish the open algebraic structure of the theory from first principles. By explicitly computing the commutator of two infinitesimal gauge transformations acting on the scalar and metric fields, we provide definitive proof that the algebra is indeed open. This calculation allows us to precisely identify the field-dependent structure functions and the coefficients multiplying the equation-of-motion terms. With the open algebra's properties firmly established, we then undertake the systematic construction of the complete BV-extended action. This process involves enlarging the field space to include a hierarchy of ghosts and antifields, whose interactions are meticulously defined by the open algebra structure. The result is an extended action, $S_{ext}$, that satisfies the classical master equation, $\{S_{ext}, S_{ext}\} = 0$, the cornerstone of the BV formalism which ensures gauge invariance is correctly encoded.The true predictive power of this framework is unleashed at the quantum level. The classical master equation is elevated to a Quantum Master Equation, which imposes a stringent consistency condition on the form of permissible quantum corrections. We leverage this equation as a powerful algebraic tool to analyze the viability of one-loop counter-terms, focusing on higher-derivative curvature invariants such as the Gauss-Bonnet scalar and the square of the Weyl tensor. Rather than attempting a direct loop calculation, our method tests which counter-terms can be consistently incorporated into the theory without violating the Quantum Master Equation. This analysis reveals a remarkable selection principle rooted in the classical open algebra. We find that counter-terms constructed from the square of the Weyl tensor are strictly forbidden, whereas a Gauss-Bonnet invariant is uniquely selected. Most strikingly, quantum consistency further constrains the coupling between the scalar field and the Gauss-Bonnet term, fixing it to be a specific logarithmic function whose parameters are inherited directly from the classical gauge transformation. This work thus uncovers a profound and predictive connection between a theory's classical algebraic structure and the definitive form of its quantum corrections, transforming the complexity of an open gauge algebra from a quantization hurdle into a powerful guiding principle.

\end{document}
                